(1) 记 \(B(r)=|E_0|(r/w)^{|\ell|}e^{-r^2/w^2}\)，选择两束初始相位为零，则 \(I(r,\phi)\propto A^2+B(r)^2+2AB(r)\cos(\ell\phi-qx)\)。若相位常数不同，只需在余弦内加入固定相位，不改变条纹拓扑。
(2) 亮条纹满足 \(\ell\phi-qx=2\pi n\)。围绕轴心一周，涡旋相位变化 \(2\pi\ell\)，因此倾斜平面波产生的条纹在轴心发生 \(|\ell|\) 条的叉形分叉；分叉朝向随 \(\ell\) 的符号反转。\(q\) 控制远离轴心的条纹间距 \(2\pi/|q|\)。
(3) 孤立涡旋强度 \(I_v\propto(r/w)^{2|\ell|}e^{-2r^2/w^2}\)。对 \(r\) 求导或对数求导，\(2|\ell|/r-4r/w^2=0\)，故 \(r_{\max}=w\sqrt{|\ell|/2}\)，而轴心强度因相位奇点为零。
(4) \(\ell\to-\ell\) 时径向强度不变而叉形方向与轨道角动量符号反转；\(A\to0\) 时干涉项消失，无法由单束强度判断符号。理想纯傍轴涡旋模每光子轴向轨道角动量为 \(\ell\hbar\)，不应与偏振自旋角动量混同。
